% !TeX root = ../main.tex
%TC:envir tabular [] word
\chapter{Conclusion}
In this project, I developed \textit{Coherence}, an actor-based programming language that combines ideas from Pony, Kappa and Lockguard to create a novel language suited for writing expressive and safe concurrent code.

This chapter summarizes the work completed, discusses potential directions for future work, and reflects on the progression of the project.

\section{Work completed}
I successfully achieved the core objectives established in the project proposal:
\begin{itemize}
  \item \textbf{Type-system}: Successfully devised a type system and implemented a type checker that enforces that all accepted programs are data-race and deadlock free. 
  \item \textbf{Compiler}: Implemented a compiler that translates \textit{Coherence} programs to LLVM IR, along with a runtime to manage scheduling and actor state.
  \item \textbf{Core Evaluation}: Developed a comprehensive test suite containing both unit and end-to-end tests to empirically verify the correctness of the implementation. Evaluated \textit{Coherence}'s expressive power by implementing representative concurrent programs that demonstrate the language's expressive power and ergonomics.
\end{itemize}
Beyond the core objectives, I pursued several extensions:
\begin{itemize}
  \item \textbf{Multi-level Pointers}: Extended the core \textit{Coherence} type system to support nested pointer structures. This involved careful translation of the viewpoint adaptation rules of Pony so that they work with \textit{Coherence}'s type system.
  \item \textbf{Performance}: Performed time-complexity analysis and ran a benchmark to evaluate the implementation of the runtime.
  \item \textbf{Formal Verification}: Wrote a Prolog program to mechanically verify that the $\triangleright$ operator satisfies properties required for soundness.
\end{itemize}
In its present form, \textit{Coherence} successfully incorporates locks and atomic sections into a reference-capability based actor model, supporting a wider range of concurrent algorithms without sacrificing data-race and deadlock freedom.

\section{Future work}
Several directions for future work remain unexplored due to time constraints and their tangential relationship to the success criteria:
\begin{itemize}
  \item \textbf{Richer language features}: 
  Introducing parametric polymorphism via generics and first-class functions would enable cleaner representation of common programming patterns. Generics could be implemented using a compile-time instantiation approach similar to C++ templates. Type-checking rules to ensure the safe use of first-class functions could also be incorporated, drawing inspiration from Pony which already support first-class functions within a capability-based type system.
  \item \textbf{Dynamic lock allocation}: I believe \textit{Coherence} could be extended to support dynamic lock allocation following the work of existing type systems\cite{dynamic_locks} while preserving deadlock freedom guarantees. This would significantly increase \textit{Coherence}'s expressive power. 
  \item \textbf{Garbage collector}: An extension that was abandoned was implementing a garbage collector for the language, as it was deemed to be tangential to the success criteria. Implementing a garbage collector would allow \textit{Coherence} to be used in long-running programs where memory efficiency is important.
\end{itemize}

\section{Reflections}
This project began with a review of literature on actor models and Pony's reference capabilities. Analysing Pony's coordination patterns revealed that certain patterns were cumbersome to express. I believed this could be solved by introducing locks. This idea became the foundation of \textit{Coherence}. Further research, drawing on work by Castegren and Gudka, refined this idea into the design of \textit{Coherence}.

Several lessons emerged from this experience. I gained an appreciation for upfront design and systematic planning. Decisions made early in the project had far-reaching consequences, and a more deliberate approach to design would have avoided some later refactoring. I also learned the value of thoroughly understanding the intent behind programming constructs, not just their mechanics. This was essential when trying to adapt Pony's viewpoint adaptation rules so that they work with \textit{Coherence}'s memory model.

Overall, this project was a success. I achieved all core objectives and several extensions, successfully combining ideas from prior work into a safe and expressive language for concurrent programming.